\section{Requirements Analysis and Scope Definition}
\label{sec:requirements-analysis}

% ---------------------------------------------------------------------------------------------------------------------
\subsection{General Requirements}
% -> es gibt zwei hautp unterschieden die man treffen kann: procativ system und reactive systeme. In are we on
% track wurden zwei verscheidene system design (aktiv und passiv) beiee hatte vor und nachteile. (Welche)
% Sind aber beide relevant.
% Daher Plattform, welche beides kann, agenda live verfolgen und agent kann selbst einschreiten.

% Im allgmeien gibt es zwei ansaätze ein solches system umzusetzen:
% 1) Intergration im bestehende Meeting-Tools wie z.B ms Zoom, Teams etc
% 2) Protoyp neu entwicklen mit einger video conference funtkionalität

% variante ein hat den vorteil, dass sich ledeglich auf das design und die implementeirung des KI-Agent
% fokussiert werden muss.allerdings ist mn+an limiert durch die integrationsmöglichkeiten, welch die planfromen
% über die API bereitstellen. Dies varriert je nach platrfom. es ist erschwer mölgich anpasseungen an dem
% vorgeben UI der Konferenzen forzuhenhemn oder sogar unmöglich. somit ist der läsungsruam  f+r den prototyp limiert.

% Variante 2) Vorteil komplett frei, nachtiel user nicht vertraun mti interfa

% Das vollständige Anforderungsbild eines idealen AI-Facilitators.
% Hier weiterhin die Unterscheidung funktional / nicht-funktional unterbringen (z.B. als kurze Absätze oder ein
% \paragraph{Functional} / \paragraph{Non-functional}), plus deine Punkte proaktiv-vs-reaktiv und Plattform-Optionen.

Der Agent sollt enicht zu schnell versuchen meinungsverscheindeheiten aufzulösen. (vgl hidden profile problem)

% ---------------------------------------------------------------------------------------------------------------------s
\subsection{Scoped Requirements Catalog for the Prototype (User Stories)}

% Die bewusste Priorisierung/Reduktion auf das konkret umgesetzte Teilproblem (Hidden Profiles /
% Shared-Information-Bias, ungleiche Beteiligung).
% Mit expliziter Begründung: Machbarkeit + Evaluierbarkeit (Anknüpfung an den Kapitel-Intro).
% Was bewusst nicht umgesetzt wird → Verweis auf Limitations/Future Work.
% Alternative Titel für 4.1.2, falls dir „Scoped" zu wiederholt ist: Prioritized and Specified Requirements
% oder Derived Prototype Requirements and Delimitation.

Die Anforderungen werden in Form von User Stories festgehalten. Laut Ott können diese
einen in Alltagssprache formulierten Software-Anforderungskatalog darstellen
\cite[p.~18]{ottAppsEffektivManagen2019}. Die Formulierungen orientieren sich an diesem Muster:
\begin{quote}
  \enquote{Als <Rolle>, ich möchte <Ziel/Wunsch>, um <Nutzen>.}
  \cite[p.~18]{ottAppsEffektivManagen2019}
\end{quote}

Da die Anforderungen aus der Literatur und den Experteninterviews extrahiert wurden, sind die User Stories
aus der Perspektive eines AI-Facilitators geschrieben, der zu einem möglichst guten Meeting beitragen möchte.
\autoref{tab:ai-facilitator-user-stories} bündelt acht zentrale Stories. Sie beschreiben das wünschenswerte
Verhaltensrepertoire, nicht bereits den implementierten Umfang des Prototyps.

\begin{table}[H]
  \centering
  \caption{User stories for an AI facilitator, derived from the facilitation literature and the expert
  interviews.}
  \label{tab:ai-facilitator-user-stories}
  \footnotesize
  \setlength{\tabcolsep}{3.5pt}
  \begin{tabularx}{\textwidth}{@{}
      >{\raggedright\arraybackslash}X
      >{\raggedright\arraybackslash}p{3.35cm}
      >{\raggedright\arraybackslash}p{3.45cm}
    @{}}
    \toprule
    \textbf{User Story} & \textbf{Reasoning} & \textbf{Source} \\
    \midrule
    Als AI-Facilitator möchte ich Ablauf und Ziel des Meetings klar machen können, damit die Gruppe
    einen gemeinsamen Bezugsrahmen hat und zu einer tragfähigen Entscheidung kommt.
    & Kernaufgabe der Facilitation ist das Strukturieren von Aufgaben und das Erreichen der
    Meeting-Outcomes. Ohne kommuniziertes Ziel und Prozess verliert die Diskussion den Fokus.
    & Hayne \cite[p.~73]{hayneFacilitatorsPerspectiveMeetings};
    Schwarz \cite{schwarzSkilledFacilitatorComprehensive2016};
    Doodle \cite[p.~11]{DoodleStateMeetings2019};
    Chen et~al.\ \cite{chenAreWeTrack2025};
    Interviews
    (\autoref{ch:empirical-analysis-of-facilitating-strategies}) \\
    \addlinespace
    Als AI-Facilitator möchte ich Personen, die noch wenig gesagt haben, auf rücksichtsvolle Weise in
    die Diskussion einladen, damit alle relevanten Stimmen gehört werden, ohne Schweigen zu
    stigmatisieren.
    & Ungleiche Beteiligung ist ein zentrales Meeting-Problem. Expert:innen laden stille Personen
    aktiv ein, warnen im Remote-Setting aber vor Bloßstellung und einem Kontrollgefühl.
    & Atlassian \cite{atlassianWorkplaceWoesMeetings};
    Chen et~al.\ \cite{chenAreWeTrack2025};
    Doodle \cite[p.~6]{DoodleStateMeetings2019};
    Interviews
    (\autoref{ch:empirical-analysis-of-facilitating-strategies}) \\
    \addlinespace
    Als AI-Facilitator möchte ich gezielt nach noch nicht genannten oder nicht ins Bild passenden
    Informationen fragen, damit ungeteiltes Wissen in die Entscheidung einfließen kann.
    & Hidden Profiles entstehen, weil Gruppen vorwiegend geteilte Informationen diskutieren.
    Nachhaken nach abweichenden Informationen ist eine belegte Gegenstrategie.
    & Stasser und Titus \cite{stasserExpertRolesInformation1995};
    Baker \cite{bakerEnhancingGroupDecision2010};
    Arkes et~al.\ \cite{arkesInformationExchangeGroup2009};
    Interviews
    (\autoref{ch:empirical-analysis-of-facilitating-strategies}) \\
    \addlinespace
    Als AI-Facilitator möchte ich Meinungsverschiedenheiten nicht vorschnell auflösen, damit Dissent
    den Informationsaustausch vertieft und bessere Alternativen sichtbar bleiben.
    & Dissent erhöht die Diskussionsintensität und die Lösungsrate von Hidden Profiles. Ein zu
    frühes Glätten würde genau diesen Effekt unterbinden (Kaners Groan Zone).
    & Schulz-Hardt et~al.\ \cite{schulz-hardtGroupDecisionMaking2006};
    Kaner \cite{kanerFacilitatorsGuideParticipatory2014};
    Interviews
    (\autoref{ch:empirical-analysis-of-facilitating-strategies}) \\
    \addlinespace
    Als AI-Facilitator möchte ich nur sparsam und mit klarem Prozessnutzen eingreifen, damit die
    menschliche Diskussion nicht übersteuert wird.
    & Facilitation ist substantiell neutral und führt die Gruppe, ohne sie zu dominieren. Zu häufige
    Eingriffe erzeugen Misstrauen; Expert:innen warnen vor Over-Facilitation und fehlender Präsenz.
    & Schwarz \cite{schwarzSkilledFacilitatorComprehensive2016};
    Houtti et~al.\ \cite{houttiObserveAskIntervene2025};
    Interviews
    (\autoref{ch:empirical-analysis-of-facilitating-strategies}) \\
    \addlinespace
    Als AI-Facilitator möchte ich thematische Abweichungen und Zeitdruck erkennen und die Gruppe
    zurück zur Agenda führen, damit das Meeting rechtzeitig zu einem Abschluss kommt.
    & Fehlender Fokus, Abschweifen und ausbleibende Entscheidungen sind klassische GSS-Probleme.
    Timekeeping und Agenda-Führung wurden in den Interviews als besonders geeigneter KI-Beitrag genannt.
    & Beranek et~al.\ \cite[p.~152]{beranekFacilitationGroupSupport1993};
    Kauffeld und Lehmann-Willenbrock \cite{kauffeldMeetingsMatterEffects2012};
    Chen et~al.\ \cite{chenAreWeTrack2025};
    Interviews
    (\autoref{ch:empirical-analysis-of-facilitating-strategies}) \\
    \addlinespace
    Als AI-Facilitator möchte ich den Diskussionsstand zusammenfassen und Entscheidungsoptionen
    sichtbar machen, damit die Gruppe aus der entstehenden Komplexität nächste Schritte ableiten kann.
    & Zusammenführen, Clustern und ein Metablick in der konvergenten Phase sind Kernpraktiken.
    Expert:innen sehen in schnellem Sensemaking den stärksten konkreten Mehrwert von KI.
    & Ackermann \cite[p.~629]{ackermannGroupSupportSystems2021};
    Kaner \cite{kanerFacilitatorsGuideParticipatory2014};
    Park et~al.\ \cite{parkCoExplorerTechnologyProbe2024};
    Interviews
    (\autoref{ch:empirical-analysis-of-facilitating-strategies}) \\
    \addlinespace
    Als AI-Facilitator möchte ich inhaltlich neutral bleiben und selbst keine Entscheidung treffen,
    damit die Gruppe die Verantwortung für das Ergebnis behält.
    & Der Facilitator ist akzeptiert, substantiell neutral und ohne Entscheidungsautorität. Die
    Expert:innen laden Leadership zur Einordnung ein, entscheiden aber nicht an Stelle der Gruppe.
    & Schwarz \cite{schwarzSkilledFacilitatorComprehensive2016};
    Hayne \cite[p.~73]{hayneFacilitatorsPerspectiveMeetings};
    Interviews
    (\autoref{ch:empirical-analysis-of-facilitating-strategies}) \\
    \bottomrule
  \end{tabularx}
\end{table}
\FloatBarrier

Für den Prototyp werden daraus bewusst nur die Stories priorisiert, die Hidden Profiles und ungleiche
Beteiligung betreffen: das Einladen stiller Stimmen, das Heben ungeteilter Informationen, das Aushalten von
Dissent, sparsame Interventionen und inhaltliche Neutralität. Agenda-Führung und ausführliches Sensemaking
bleiben wünschenswert, liegen aber außerhalb des evaluierten Interventionsmechanismus.
